\documentclass{article}
\usepackage{mathnotes}

\title{Formality and Proofs}

\begin{document}

\begin{definition}[Formal System]
A  \textbf{formal system} $\mathcal{F}$ is an ordered quadruple

\[ \mathcal{F} = (\Sigma, W, A, R) \]

where:

\begin{enumerate}
\item $\Sigma$ is a finite, nonempty \textbf{alphabet} (a set of symbols);
\item $W \subseteq \Sigma^{*}$ is the set of \textbf{well-formed-formulas} (wffs), defined inductively by formation rules specifying which finite strings over $\Sigma$ belong to $W$
\item $A \subseteq W$ is a distinguished subset of $W$, whose members are called \textbf{axioms.}
\item $R$ is a finite set of \textbf{rules of inference}, each of which is a relation on $W$ determining from which formulas other formulas may be derived.
\end{enumerate}
\end{definition}

\begin{definition}[Derivation]\synonyms{proof}
A \textbf{derivation} (or \textbf{proof}) in $\mathcal{F}$ is a finite sequence of formulas

\[ \varphi_1, \varphi_2, \dots, \varphi_n \]

such that for each $i \le n$,

\begin{enumerate}
\item Either $\varphi_i \in A$, or
\item $\varphi_i$ follows from earlier formulas in the sequence by an inference rule in $R$.
\end{enumerate}
\end{definition}

\begin{definition}[theorem]
A formula $\psi \in W$ is a \textbf{theorem} (or \textbf{provable formula}) of $\mathcal{F}$ if there exists such a derivation ending with $\psi$. We then write

\[ \vdash_{\mathcal{F}} \psi. \]
\end{definition}

\begin{definition}[recursive]
The system $\mathcal{F}$ is said to be \textbf{recursive} (or \textbf{effectively generated}) if $W$, $A$, and $R$ are all \dref{recursively-enumerable} sets, so that derivations can be verified by a finite mechanical procedure.
\end{definition}

\begin{remark}
No \dref{semantics} are assumed in this definition; a \dref{formal-system} is purely \dref{syntactic}.

A \dref{formal-system} operates purely on \dref{form}, not content.

Its symbols are inert marks; its rules manipulate these marks according to syntactic shape alone.

No step depends on what any symbol "means."

A proof in such a system is therefore a finite, mechanical transformation of strings; a sequence of moves in a symbol game governed entirely by explicit rules.

Semantics enter only afterward, when we interpret the symbols externally (e.g. as numbers, sets, or truth values).

Until then, the system is a syntax engine; a rule-bound procedure that generates certain strings (theorems) from others (axioms), without regard to interpretation or truth.
\end{remark}

\begin{note}
When an interpretation (model) $\mathcal{M}$ and a satisfaction relation $\models_{\mathcal{M}}$ are supplied, one obtains a \textbf{formal theory.}
\end{note}

\end{document}
